% ======================================================================
% Ontologie der Kontinua -- Core 1.3.3
% Cross-K-Modul: crossk_k3_k4.tex
% Cross-Level-Relationen zwischen K3 und K4
% ======================================================================

\subsubsection{K3–K4-Querverbindung}
\label{sec:crossk-k3-k4}

Der Übergang $K_3 \rightarrow K_4$ markiert die Entstehung des biologischen
Kontinuums aus einem chemischen Kontinuum, das bereits Geometrie, Reaktionen und
dynamische Flüsse besitzt. Während $K_3$ chemische Konnektivität, Reaktionspfade
und stabile Cluster bereitstellt, führt $K_4$ ein:
\begin{itemize}
  \item Membranen und Randbildung,
  \item andauernde Gradienten (Ionen-, pH-, Konzentrations- und Redox-Gradienten),
  \item osmotische und elektrochemische Potentiale,
  \item aktive und passive Transportflüsse,
  \item metabolische Zyklen und Pufferungsschleifen,
  \item frühe Informationsträger (Ladung, Struktur, Templating),
  \item regulatorische Schwellen und proto-homeostatische Bedingungen.
\end{itemize}

Dies ist die Geburt der ersten lebensähnlichen Kontinua: Systeme, die
nichtgleichgewichtige Strukturen aufrechterhalten können.

% ----------------------------------------------------------------------
\subsubsection{1. Stufen und ihre Rollen}

\paragraph{K3 (chemisches Kontinuum).}
$K_3$ liefert:
\begin{itemize}
  \item Achsen der Zusammensetzung (Spezies, Ladung, Stöchiometrie),
  \item Potentiale aus Reaktionslandschaften,
  \item Reaktions- und Ladungsfluss-Flüsse $J_3$,
  \item stabile Cluster und chemische Zyklen,
  \item keine dauerhaften Grenzen und keine regulierten Gradienten.
\end{itemize}

\paragraph{K4 (biologisches Kontinuum).}
$K_4$ führt ein:
\begin{itemize}
  \item expliziten Rand $\partial\Omega(K_4)$: Membranen und Kompartimente,
  \item Gradienten als Achsen ($A_{\mathrm{grad}}$),
  \item Potentiale $P_{\mathrm{grad}}, P_{\mathrm{ion}}, P_{\mathrm{redox}}$,
  \item Flüsse $J_4$ (osmose-, Ionenleit-, aktive Transport- und
        metabolische Flüsse),
  \item Zyklen $C_4$ (Energie-, Pufferungs-, Redox- und Pumpzyklen),
  \item Schwellen $\Theta_4$ für Membranstabilität, Gradientenhaltung und osmotisches
        Gleichgewicht.
\end{itemize}

Die Rolle des Übergangs K₃→K₄ ist die Stabilisierung von Struktur durch
kontrollierte Nichtgleichgewichts-Dynamik.

% ----------------------------------------------------------------------
\subsubsection{2. Gemeinsame und vererbte Achsen sowie Schwellen}

\paragraph{Achsen.}
Aus $K_3$ übernimmt $K_4$:
\[
A_{\mathrm{chem}}(K_3) \subset A(K_4).
\]
Neue Achsen sind:
\begin{itemize}
  \item $A_{\mathrm{grad}}$ — Konzentrations-, pH-, Redox- und elektrische
        Gradienten,
  \item $A_{\mathrm{mem}}$ — Membrankrümmung und Permeabilitätszustände,
  \item $A_{\mathrm{exc}}$ — proto-elektrische Erregungsachse, die an K₅ anschließt.
\end{itemize}

\paragraph{Schwellen.}
$K_4$ ergänzt Schwellenfamilien:
\begin{itemize}
  \item $\Theta_{\mathrm{mem}}$ — Membranriss-/Permeabilitätsgrenzen,
  \item $\Theta_{\mathrm{grad}}$ — minimale und maximale tragfähige Gradienten,
  \item $\Theta_{\mathrm{osm}}$ — osmotische Druckeinschränkungen ($\Delta\pi$),
  \item $\Theta_{\mathrm{redox}}$ — Stabilität der Redox-Bilanz,
  \item $\Theta_{\mathrm{energy}}$ — Schwellen metabolischer Lebensfähigkeit.
\end{itemize}

Wenn $\Theta_3$ versagt (chemische Stabilität), bildet sich keine biologische
Grenze:
\[
\Theta_3^{\mathrm{death}} \Rightarrow \Omega(K_4)=\varnothing.
\]

% ----------------------------------------------------------------------
\subsubsection{3. Cross-Level-Flüsse, Zyklen und Spannungen}

\paragraph{Flüsse.}
Flüsse $J_4$ schließen alle chemischen Flüsse von $J_3$ ein plus:
\begin{itemize}
  \item $J_{\mathrm{osm}}$ — osmotischer Wasserfluss,
  \item $J_{\mathrm{ion}}$ — Ionenfluss über Kanäle oder Membran,
  \item $J_{\mathrm{pump}}$ — aktiver Transport gesteuert durch $P_{\mathrm{energy}}$,
  \item $J_{\mathrm{redox}}$ — Elektronentransfer und Energieumwandlung,
  \item $J_{\mathrm{metabolic}}$ — gekoppelte Reaktions–Transport-Schleifen.
\end{itemize}

Diese definieren das erste nichtgleichgewichtige Dynamikregime der Hierarchie.

\paragraph{Zyklen.}
Zyklen $C_4$ (Energie, Pufferung, Pumpen, Redox) sind strukturelle Stabilisatoren:
\[
C_{\mathrm{energy}} : P_{\mathrm{source}} \rightarrow J_{\mathrm{pump}}
\rightarrow \text{Gradientenwiederherstellung} \rightarrow P_{\mathrm{source}}.
\]

Ein Zyklus besteht nur, wenn:
\[
\oint_C dA > 0 \quad\text{und}\quad S(C) > 0,
\]
wobei $S(C)$ die Zyklusstabilität aus der Core-Definition ist.

\paragraph{Spannung.}
Die Cross-Level-Spannung $T_{3,4}$ reflektiert Inkongruenzen zwischen chemischer
Reaktivität und biologischer Stabilität:
\[
T_{3,4} = f(\Theta_{\mathrm{mem}} - P_{\mathrm{chem}},\
           \Theta_{\mathrm{grad}} - P_{\mathrm{grad}},\
           J_{\mathrm{react}} - J_{\mathrm{buffer}}).
\]
Übermäßige Spannung zerstört Gradienten oder Membranen und kollabiert $K_4$
zurück zu $K_3$.

% ----------------------------------------------------------------------
\subsubsection{4. Geburts- und Todesbedingungen über die Stufen}

\paragraph{Geburt von \texorpdfstring{$K_4$}{K_4}.}
Das biologische Kontinuum tritt auf, wenn:
\[
T_3 > \Theta_{3,\mathrm{dim}}
\quad \text{und} \quad
A_{\mathrm{grad}}, A_{\mathrm{mem}} \in M_4 \setminus A(K_3),
\]
was der Entstehung von Grenzen und persistenten Gradienten entspricht.

Der Zustandsraum erweitert sich als:
\[
\Omega(K_4) = \Omega(K_3) \cup \Delta\Omega_{\mathrm{bio}}.
\]

\paragraph{Todesausbreitung.}
Wenn Membranen keine Gradienten halten können (Verletzung von
$\Theta_{\mathrm{mem}}$, $\Theta_{\mathrm{grad}}$, $\Theta_{\mathrm{osm}}$), gilt:
\[
\Omega(K_4) \rightarrow \Omega(K_3),
\]
also ein Kollaps in das chemische Regime.

$K_3$ bleibt stabil, solange dessen eigene Schwellen nicht brechen.

% ----------------------------------------------------------------------
\subsubsection{5. Konzeptionelle Beispiele}

\paragraph{Osmotische Grenzbildung.}
Die Membran entsteht als stabiler Bereich von $\partial\Omega$, der erfüllt:
\[
\Delta\pi = RT(C_{\mathrm{in}} - C_{\mathrm{out}})
\]
innerhalb von $\Theta_{\mathrm{osm}}$.

\paragraph{Gradientenbasierte Organisation.}
Ein stabiler $pH$- oder Iongradient definiert eine neue Achse
$A_{\mathrm{grad}}$, die ermöglicht:
\begin{itemize}
  \item Proto-Metabolismus,
  \item gerichtete Energieumwandlung,
  \item selektiven Transport.
\end{itemize}

\paragraph{Proto-Erregung.}
Elektrische Gradienten $\Delta V$ ermöglichen proto-spike-Verhalten, wenn:
\[
\Delta V > \Theta_{\mathrm{exc}},
\]
was $K₄$ an K₅ über frühe informationsartige Flüsse bindet.

Diese Beispiele veranschaulichen die Essenz des K3→K4-Übergangs in komprimierter
Form: Grenzen, Gradienten und nichtgleichgewichtige Zyklen werden zur ersten
stabilen biologischen Organisation.
